\section{Introduction}\label{sec:intro}

An urban traffic engineer retiming signalized intersections during an unexpected demand surge faces two coupled statistical hurdles: identifying low-delay phase plans across dozens of signal phases under a tight simulation budget, and certifying before field deployment that the chosen plan satisfies an emissions or congestion chance constraint with high confidence. Similar decision problems arise when supply-chain managers calibrate dual-sourcing replenishment rules, service operators schedule time-varying staffing and surge diversion, and power-grid operators dispatch battery reserves against volatile net load \citep{boesel2003using,ankenman2010stochastic,hong2011sequential,ault2021reinforcement}. Across these settings, policy vectors span $d_t \in [44, 336]$ continuous decision variables, simulation noise is strongly heteroscedastic across operating states \citep{binois2018practical}, and noisy search evaluations cannot be reused naively to certify safety without optimistic selection bias. Consequently, prudent deployment separates sequential search from an independent terminal clean-up test and defaults to a pre-existing status-quo policy upon abstention—even though under an exogenous demand surge, that status-quo policy may itself violate the chance constraint and require active compliance restoration.

Organizations rarely optimize each planning problem from scratch. Enterprise simulation repositories typically hold historical solver trajectories (\texttt{SPSA}, \texttt{CMA-ES}, evolutionary heuristics) logged on shorter planning horizons or smaller subnetworks ($d_s \neq d_t$). Yet turning such ``dark-data'' logs into a safe, dimension-invariant prior for chance-constrained simulation optimization raises four fundamental challenges:
\begin{enumerate}
  \item \emph{Dimensional mismatch and adaptive archive bias.} Source and target policies live in different Euclidean spaces ($[0,1]^{d_s}$ vs.\ $[0,1]^{d_t}$), and legacy optimizers heavily oversample narrow elite basins while logging many single-replication screening evaluations ($r_a = 1$) alongside multi-replication runs \citep{spall1992multivariate,hansen2006cma,deb2002fast}.
  \item \emph{Mean-shift leakage into heteroscedastic variance.} When a target task experiences a structural regime shift relative to historical tasks, fitting noise models to residuals around an inherited source mean conflates deterministic prior mismatch with intrinsic Monte Carlo variance, inflating uncertainty estimates and distorting feasibility boundaries.
  \item \emph{Unbounded shift risk and finite-source covariance instability.} With a modest number of historical source tasks ($M_{\mathrm{src}} = 6$) relative to basis dimension, empirical between-task covariances are rank-deficient and sensitive to individual source tasks, while fixed-weight transfer priors suffer severe asymptotic drag if prior--target discrepancies are large \citep{grunwald2017inconsistency,wang2024pretrained}.
  \item \emph{Misalignment between search acquisition and multi-slot clean-up certification.} Standard constrained Bayesian optimization scores queries by marginal probability of feasibility $\Prob(g_t(x) \le 0) \ge 1 - \alpha$ \citep{eriksson2021scalable,ungredda2024bayesian}. Under an exact binomial verification gate of $K_0 = \SlotZeroK$ replications at level $\delta_0 = \SlotZeroDelta$ and $\alpha = 0.05$, however, a policy sitting exactly on the $5\%$ risk boundary passes with probability only $C_{\SlotZeroK,1}(0.05) \approx 0.0221$. Moreover, when terminal clean-up tests a guarded three-slot shortlist over already-evaluated policies $\mathcal{D}_n$, querying a candidate $q$ alters future deployment loss through two distinct channels: refining posterior beliefs over $\mathcal{D}_n$ and admitting an unevaluated $q \notin \mathcal{D}_n$ into verification eligibility $\mathcal{D}_n \cup \{q\}$.
\end{enumerate}

\subsection{Summary of Contributions}

We develop and analyze \emph{Gate-Power-Aligned Functional Transfer} (\textbf{GP-FTA}), closing each of the four challenges above within a unified statistical and algorithmic architecture:
\begin{enumerate}
  \item \textbf{Public Exposure Rollout, Occupancy-Balanced Basis, and Shift-Invariant Variance (Section~\ref{sec:prior}).} Deterministic rollouts driven solely by public planning forecasts map policies of arbitrary horizon $d_k$ onto a $\PsiDim$-dimensional state-risk exposure $\psi(x)$ and volatility descriptors $(A_{\mathrm{local}}(x), N_{\mathrm{shared}}(x))$. Clipped inverse-occupancy weighting, graph-Laplacian low-pass filtering, and Gram whitening extract $K_{\mathrm{act}} = \ActiveRank$ orthonormal functional modes while routing pruned spectral energy into online GP conditioning features. Heteroscedastic variance is fitted via nonnegative least squares exclusively to within-policy sums of squares $S^2/(r-1)$ on replicated runs ($r \ge 2$), preserving exact algebraic invariance to arbitrary mean shifts while retaining $r_a = 1$ archive records for mean estimation. Proposition~\ref{prop:subspace} and Theorem~\ref{thm:transfer_boundary} establish self-normalized clipping bias bounds and the exact scalar signal-to-mismatch boundary for transfer dominance.
  \item \textbf{Design-Whitened Fisher-Eigenspace \texttt{SURE} Shrinkage and Covariance Stabilization (Sections~\ref{sec:model} and~\ref{sec:theory}).} Projecting centered target observations onto the eigenspace of the target Fisher information matrix yields unit-variance, pairwise-orthogonal subspace coordinates on which Stein--Mallows unbiased risk estimation (\texttt{SURE}) governs smooth shrinkage and source-prior release, paired with dimension-aware finite-source covariance shrinkage $\omega_{\mathrm{src}} = p/(M_{\mathrm{src}}+p)$. Theorem~\ref{thm:sure_uniform} proves a deterministic self-normalized excess-loss bound over target-only release that holds \emph{uniformly over all prior--target shifts} $\|\bm{\delta}\|_2 \in (0, \infty)$, eliminating grid-cardinality union-bound slack.
  \item \textbf{Eligible-Cascade Knowledge Gradient and Guarded Exact-Binomial Clean-Up (Sections~\ref{sec:vkg} and~\ref{sec:theory}).} We formulate a fast two-stage \emph{Eligible-Cascade Knowledge Gradient} (\textbf{EC-KG}) acquisition rule over the post-query eligible set $\mathcal{D}_n \cup \{q\}$ and prove an exact additive decomposition $\mathrm{VOI}_n^{\mathrm{EC}}(q) = \mathrm{VOI}_n^{\mathrm{refine}}(q) + \mathrm{VOI}_n^{\mathrm{admit}}(q) \ge 0$ (Theorem~\ref{thm:eckg_dual_voi}), where $\mathrm{VOI}_n^{\mathrm{admit}}(q) = 0$ on revisits $q \in \mathcal{D}_n$. At termination, an objective-guarded three-slot exact-binomial clean-up cascade guarantees finite-sample familywise false-certification control $\le \delta$ under the verification law without distributional assumptions (Theorem~\ref{thm:vkg_safety}). Key theoretical identities are mechanized in Lean~4.31.0 with zero \texttt{sorry} axioms.
  \item \textbf{Prospective Factorial Confirmation, 96-State Census, and Cross-Domain Validation (Sections~\ref{sec:experiments}--\ref{sec:results}).} In a pre-registered 720-run prospective factorial study across 72 unseen inventory, queueing, and surge-traffic tasks, \textbf{GP-FTA-ECKG} achieves 130/130 confidently safe certifications (zero empirical false certifications), significantly outperforms business-feasibility KG across the 72-task panel (Holm-adjusted two-sided $p = 0.039$; $p = 0.027$ in traffic with a $9.03$-percentage-point reduction in operating violation rate), wins $44$ of $72$ paired tasks against single-slot gate-power KG (Wilcoxon $p = 0.014$), and beats cold-start public-representation Thompson sampling in inventory ($p = 0.029$) and queueing (Holm $p = 0.002$) while releasing cleanly under severe traffic surge mismatch. A complete 96-state early/near-terminal query census and a 9,796-run four-family benchmark—spanning ENTSO-E grid storage, single-replication enterprise archives, and microscopic \texttt{Ingolstadt21} SUMO—delineate both the operational gains and the structural boundaries of legacy transfer.
\end{enumerate}
